\documentclass[11pt,a4paper]{article}
\usepackage[utf8]{inputenc}
\usepackage[T1]{fontenc}
\usepackage[french]{babel}
\usepackage{amsmath,amssymb,amsthm}
\usepackage{amsfonts}
\usepackage{mathtools}
\usepackage{geometry}
\usepackage{hyperref}
\usepackage{enumitem}
\usepackage{booktabs}
\usepackage{array}
\usepackage{mathrsfs}
\usepackage{xcolor}
\usepackage{microtype}
\usepackage{tikz-cd}
\usepackage{graphicx}
\usepackage{stmaryrd}
\usepackage{upgreek}

\geometry{a4paper, top=2.5cm, bottom=2.5cm, left=2.5cm, right=2.5cm}

\theoremstyle{plain}
\newtheorem{theorem}{Th\'eor\`eme}[section]
\newtheorem{lemma}[theorem]{Lemme}
\newtheorem{lemme}[theorem]{Lemme}
\newtheorem{proposition}[theorem]{Proposition}
\newtheorem{corollary}[theorem]{Corollaire}
\newtheorem{conjecture}[theorem]{Conjecture}
\newtheorem{theoreme}[theorem]{Th\'eor\`eme}
\newtheorem{surprise}[theorem]{Surprise}

\theoremstyle{definition}
\newtheorem{definition}[theorem]{D\'efinition}
\newtheorem{example}[theorem]{Exemple}
\newtheorem{remark}[theorem]{Remarque}
\newtheorem{variante}[theorem]{Variante}

\DeclareMathOperator{\Ext}{Ext}
\DeclareMathOperator{\Hom}{Hom}
\DeclareMathOperator{\Ker}{Ker}
\DeclareMathOperator{\Coker}{Coker}
\DeclareMathOperator{\Spec}{Spec}
\DeclareMathOperator{\Specf}{Spf}
\DeclareMathOperator{\Proj}{Proj}
\DeclareMathOperator{\Gal}{Gal}
\DeclareMathOperator{\Aut}{Aut}
\DeclareMathOperator{\End}{End}
\DeclareMathOperator{\Pic}{Pic}
\DeclareMathOperator{\Frob}{Frob}
\DeclareMathOperator{\Id}{Id}
\DeclareMathOperator{\Tr}{Tr}
\DeclareMathOperator{\rk}{rg}
\DeclareMathOperator{\detZ}{det}
\DeclareMathOperator{\Fr}{Fr}
\DeclareMathOperator{\Gri}{Gr}
\DeclareMathOperator{\diag}{diag}

\newcommand{\Q}{\mathbb{Q}}
\newcommand{\Z}{\mathbb{Z}}
\newcommand{\F}{\mathbb{F}}
\newcommand{\C}{\mathbb{C}}
\newcommand{\N}{\mathbb{N}}
\newcommand{\Ql}{\overline{\mathbb{Q}}_\ell}
\newcommand{\Qlt}{\overline{\mathbb{Q}}_{\ell'}}
\newcommand{\Zl}{\overline{\mathbb{Z}}_\ell}
\newcommand{\Fp}{\mathbb{F}_p}
\newcommand{\car}{\text{car}}

\title{\textbf{Comptage de faisceaux $\ell$-adiques}}
\author{par Pierre Deligne\\[.5em]
\small \`a G\'erard Laumon, \`a l'occasion de son soixanti\`eme anniversaire}
\date{}

\begin{document}

\maketitle

\section{Introduction}

\subsection{}
Soient $X_0$ une courbe projective, lisse, absolument connexe de genre $g$, sur un corps fini $\F_q$ de caract\'eristique $p$, et $X$ celle qui s'en d\'eduit par extension des scalaires \`a une cl\^oture alg\'ebrique $F$ de $\F_q$:
\begin{equation}
\begin{tikzcd}
X \arrow[r] \arrow[d] & X_0 \arrow[d] \\
\Spec(F) \arrow[r] & \Spec(\F_q)
\end{tikzcd}
\tag{1.1.1}
\end{equation}

L'endomorphisme de $X_0$ qui est l'identit\'e sur l'espace topologique sous-jacent, et $f \mapsto f^q$ sur le faisceau structural, est un endomorphisme de $\F_q$-sch\'ema. Nous noterons $\Frob$ l'endomorphisme du $F$-sch\'ema $X$ qui s'en d\'eduit par extension des scalaires. C'est l'endomorphisme de Frobenius de $X$. Fixons un nombre premier $\ell \neq p$, une cl\^oture alg\'ebrique $\Ql$ de $\Q_\ell$, et soit $E$ l'ensemble des classes d'isomorphie de $\Ql$-faisceaux lisses irr\'eductibles de rang $2$ sur $X$. L'image inverse par $\Frob$ induit une permutation de $E$. Nous la noterons $V$.

\subsection{}
Dans l'article [Dr], qui reste pour moi aussi myst\'erieux qu'il y a $31$ ans, Drinfeld calcule le nombre de points fixes de $V: E \to E$.

Un $\Ql$-faisceau $\mathcal{L}_0$ de rang un sur $\Spec(\F_q)$ est d\'etermin\'e \`a isomorphisme pr\`es par l'unit\'e $\lambda$ de $\Ql$ telle que le Frobenius g\'eom\'etrique $\Fr \in \Gal(F/\F_q)$ agisse par multiplication par $\lambda$ sur la fibre de $\mathcal{L}_0$ au point g\'eom\'etrique $\Spec(F)$. La $\F_q$-torsion de $F_0$ sur $X_0$ par $\mathcal{L}_0$ est le produit tensoriel avec l'image inverse de $\mathcal{L}_0$ sur $X_0$. Par abus de langage, on dira aussi ``$\F_q$-torsion par $\lambda$''.

Drinfeld utilise que la classe d'isomorphie d'un $\Ql$-faisceau lisse $F$ sur $X$ est fixe par $\Frob^*$ si et seulement si $F$ est l'image inverse d'un $\Ql$-faisceau $F_0$ sur $X_0$, et que, si $F$ est irr\'eductible, $F_0$ est unique \`a $\F_q$-torsion pr\`es. Le probl\`eme r\'esolu par [Dr] devient celui de compter les $\Ql$-faisceaux lisses de rang $2$ sur $X_0$, pris \`a $\F_q$-torsion pr\`es, et en ne consid\'erant que ceux qui sont irr\'eductibles et le restent apr\`es image inverse sur $X$.

En 1981, on disposait presque de la correspondance entre $\Ql$-faisceaux lisses irr\'eductibles de rang $2$ sur $X$ et repr\'esentations automorphes cuspidales partout non ramifi\'ees pour $GL(2,k(X_0))$. Gr\^ace \`a celle-ci, Drinfeld ramenait le probl\`eme \`a une application de la formule des traces pour $GL(2)$. Cette r\'eduction montre que le nombre cherch\'e est ind\'ependant de $\ell$. La formule obtenue par Drinfeld a les propri\'et\'es miraculeuses (A) et (B) suivantes.

Pour $n \geq 1$, soit $N_n$ le nombre de points fixes de l'it\'er\'e $n$i\`eme $V^n: E \to E$ de $V$. Calculer $N_n$ est le probl\`eme ci-dessus, avec $X_0/\F_q$ remplac\'e par $(X_0 \otimes_{\F_q} \F_{q^n})/\F_{q^n}$.

\medskip
\noindent\textbf{(A)} La fonction $n \mapsto N_n$ a la forme
\begin{equation}
N_n = \sum_i a_i \beta_i^n
\tag{1.2.1}
\end{equation}
pour des entiers $a_i$ et des nombres de Weil $\beta_i$ convenables.

Les $\beta_i$ sont des mon\^omes en $q$ et en les valeurs propres de l'endomorphisme $\Frob^*$ de $H^1(X, \Ql)$. Le genre $g$ \'etant fix\'e, quels mon\^omes $\beta_i$ apparaissent et avec quelles multiplicit\'es $a_i$ ne d\'epend pas de la courbe consid\'er\'ee.

La formule des traces exprime $N_n$ comme somme de plusieurs termes. Pris s\'epar\'ement, ces termes n'ont pas tous la forme (1.2.1).

Soit $\mathcal{P}$ une surface de Riemann compacte de genre $g$. Consid\'erons les syst\`emes locaux d'espaces vectoriels complexes sur $\mathcal{P}$. Soit $M$ l'espace de modules de ceux qui sont irr\'eductibles de rang $2$. L'espace $M$ et l'ensemble $E$ sont vides si $g = 0$ ou $1$. Sinon, $M$ est une vari\'et\'e symplectique complexe connexe. Sa dimension complexe est donc un entier pair $2N$.

\medskip
\noindent\textbf{(B)} Dans (1.2.1), le terme dominant est $q^N$: un des $\beta_i$ est $q^N$, sa multiplicit\'e $a_i$ est $1$, et les autres $\beta_i$ v\'erifient $|\beta_i| < q^N$.

\subsection{}
La formule (1.2.1) est r\'eminiscente d'une formule des traces de Lefschetz o\`u, dans de bons cas, le nombre de points fixes des it\'er\'es $T^n$ ($n \geq 1$) d'un endomorphisme $T$ d'un espace $S$ est
\begin{equation}
\Tr(T^* \mid H^*_?(S)) := \sum_i (-1)^i \Tr(T^* \mid H^i_?(S)).
\tag{1.3.1}
\end{equation}

Le ``$?$'' est l\`a pour rappeler qu'il faut consid\'erer une cohomologie avec conditions de support. Les conditions de support \`a imposer d\'ependent du comportement \`a l'infini de $T$. Par exemple, si pour une fonction d'exhaustion $f$ sur $S$ on a $f(T(x)) > f(x)$ (resp.\ $f(T(x)) < f(x)$) pour $f(x)$ assez grand, la cohomologie \`a consid\'erer est la cohomologie \`a support compact (resp.\ ordinaire).

Supposons que $H^*_?$ soit de dimension finie, pour que (1.3.1) ait un sens. Le membre de droite de (1.3.1) est alors de la forme (1.2.1): c'est $\sum a(\beta)\beta^n$, o\`u $\beta$ parcourt les valeurs propres de $T^*$ et o\`u l'entier $a(\beta)$ est la somme altern\'ee des multiplicit\'es de $\beta$ comme valeur propre des endomorphismes $T^*$ des $H^i_?(S)$.

Je n'ai malheureusement aucune id\'ee quant \`a comment consid\'erer l'ensemble $E$ des classes d'isomorphie de $\Ql$-faisceaux lisses irr\'eductibles de rang $2$ sur $X$ comme un ``espace'' ayant une cohomologie $H^*_?(E)$ telle qu'on puisse esp\'erer que le nombre de points fixes de $V^n$ soit donn\'e par une formule du type (1.3.1).

\subsection{}
M\^eme si on ne sait pas comment penser g\'eom\'etriquement \`a $E$, si $f \in E$ est la classe d'isomorphie de $F$, on sait ce que devrait \^etre le compl\'et\'e formel $E^\wedge_f$ de $E$ en $f$. Une $\Ql$-alg\`ebre locale de dimension finie $\Lambda$ est augment\'ee vers $\Ql$. Une d\'eformation de $F$ sur $\Spec(\Lambda)$ est un $\Ql$-faisceau lisse $F_\Lambda$ sur $X$, muni d'une structure de $\Lambda$-module et d'un isomorphisme $F_\Lambda \otimes_\Lambda \Ql \xrightarrow{\sim} F$, tel qu'en un point $x$ (et donc en tout point $x$ de $X$), $(F_\Lambda)_x$ soit un $\Lambda$-module libre. Le foncteur en $\Spec(\Lambda)$ des classes d'isomorphie de d\'eformations de $F$ sur $\Spec(\Lambda)$ est prorepr\'esentable. Notons $R_F$ la $\Ql$-alg\`ebre locale compl\`ete de corps r\'esiduel $\Ql$ telle que $\Specf(R_F)$ le prorepr\'esente. Il n'y a pas d'obstruction aux d\'eformations, et $R_F$ est donc isomorphe \`a une alg\'ebre de s\'eries formelles $\Ql[[t_1,\ldots,t_m]]$.

Parce que $F$ est irr\'eductible, ses seuls automorphismes sont les multiplications par un $\lambda \in \Ql^*$. Ils se prolongent \`a toute d\'eformation. Si $F$ et $F'$ sont deux repr\'esentants de la classe d'isomorphie $f$, deux isomorphismes de $F$ avec $F'$ induisent donc le m\^eme isomorphisme entre $R_F$ et $R_{F'}$: $E^\wedge_f := \Specf(R_F)$ ne d\'epend, \`a isomorphisme unique pr\`es, que de $f$.

L'espace tangent de Zariski de $E^\wedge_f$ en $f$ est $H^1(X, \End(F))$. La forme bilin\'eaire sym\'etrique $\Tr(fg)$ sur $\End(F)$ est une autodualit\'e. Elle induit sur $H^1(X, \End(F))$ une forme symplectique \`a valeurs dans $H^2(X, \Ql) = \Ql(-1)$. La m\^eme construction donne sur $M$ la structure symplectique de $M$. Le compl\'et\'e formel $E^\wedge_f$ est donc lisse de dimension paire, la m\^eme que celle de $M$.

Le foncteur d'image inverse par $\Frob: X \to X$ transforme d\'eformation de $F$ en d\'eformation de $\Frob^*F$. Il induit donc un morphisme
\[
V^\wedge_f: E^\wedge_f \longrightarrow E^\wedge_{V(f)}.
\]

L'argument qui montre que, parce que $\Frob$ induit une \'equivalence de sites \'etales, $V$ est bijectif, montre aussi que les $V^\wedge_f$ sont des isomorphismes. Autre preuve: soit $\omega$ la forme symplectique d\'efinie ci-dessus sur l'espace tangent en $f$. Si $[t] \in H^1(X, \End(F))$ est la classe du vecteur tangent $t$ en $F$, la classe de $dV^\wedge_f(t)$ dans $H^1(X, \End(\Frob^*F))$ est $\Frob^*([t])$. Pour deux vecteurs tangents $t$ et $t'$, $\omega(dV^\wedge_f(t), dV^\wedge_f(t'))$ est donc l'image inverse, par $\Frob^*$, de $\omega(t_1, t_2) \in \Ql(-1)$, identifi\'e \`a $H^2(X, \Ql)$. Parce que $\Frob: X \to X$ est de degr\'e $q$, l'endomorphisme $\Frob^*$ de $H^2(X, \Ql)$ est la multiplication par $q$, et
\[
\omega(dV^\wedge_f(t_1), dV^\wedge_f(t_2)) = q\,\omega(t_1, t_2),
\]
ce qui s'\'ecrit aussi $(dV^\wedge_f)^*\omega = q\omega$. La diff\'erentielle $dV^\wedge_f$ est donc un isomorphisme.

L'application $dV^\wedge_f$ entre espaces tangents de Zariski s'identifie \`a
\[
\Frob^*: H^1(X, \End(F)) \longrightarrow H^1(X, \End(\Frob^*F)).
\]

Si $f$ est un point fixe de $V$, $F$ est l'image inverse de $F_0$ sur $X_0$. On sait qu'on peut choisir $F_0$ pur, et $\End(F_0)$ est donc pur de poids $0$. Les valeurs propres de $dV^\wedge_f$ sur l'espace tangent de Zariski de $E^\wedge_f$ en $f$ sont donc des nombres de Weil de poids $1$. En particulier, aucune valeur propre de $dV^\wedge_f$ en $f$ n'est \'egale \`a $1$. C'est ce qui rend raisonnable qu'une formule \`a la Lefschetz (1.3.1) puisse compter les points fixes de $V^n$, tous pris avec la multiplicit\'e un.

\subsection{}
Si $Y_0$ est un sch\'ema de type fini sur $\F_q$, l'endomorphisme de Frobenius $\Frob$ de $Y := Y_0 \otimes_{\F_q} F$ induit une permutation $\Frob$ de $Y(F)$. Ce qui pr\'ec\`ede montre que la nature de $(E,V)$ est diff\'erente de celle de $(Y, \Frob)$:

\begin{enumerate}[label=(\roman*)]
\item $E$ est ``sur $\Ql$'', donc un objet de caract\'eristique $0$.
\item Les $V^\wedge_f$ sont des isomorphismes. On aimerait dire que $V$ est de ``degr\'e un''.
\item Les points fixes de $V$ sont isol\'es car les valeurs propres de $dV^\wedge_f$ en un point fixe sont des nombres de Weil de poids $1$. Ceux de $\Frob$ sont isol\'es car $d\Frob = 0$.
\item Si $E$ est de dimension $2N$, en ce sens que les $E^\wedge_f$ le sont, l'ordre de grandeur du nombre de points fixes de $V^n$ est $(q^n)^N$, plut\^ot que $(q^n)^{2N}$.
\end{enumerate}

\subsection{}
L'espace de modules $M$ est une vari\'et\'e alg\'ebrique complexe. Contrairement \`a ce qu'une analogie h\^ative pourrait laisser croire, $E$ n'est pas de fa\c{c}on naturelle l'espace des $\Ql$-points d'une vari\'et\'e alg\'ebrique sur $\Ql$. Soit $x \in X(F)$. L'ensemble $E$ s'identifie \`a l'ensemble des classes d'isomorphie de repr\'esentations continues irr\'eductibles $V$ de $\pi_1(X,x)$, de dimension $2$ sur $\Ql$. Comme coordonn\'ees sur $E$, on aimerait prendre les $\Tr(\gamma, V)$ pour $\gamma \in \pi_1(X,x)$. Dans le cas complexe, cette construction fournit sur $M$ sa structure ``de Betti'' de vari\'et\'e alg\'ebrique complexe. Ici, $\pi_1(X,x)$ \'etant un groupe profini, toute repr\'esentation $V$ admet un r\'eseau invariant, et les $\Tr(\gamma, V)$ sont \`a valeurs dans l'anneau de valuation $\Zl$ du corps valu\'e $\Ql$.

Le point de vue ``$E$ espace de repr\'esentations de $\pi_1(X,x)$'' n'aide pas non plus \`a comprendre que le nombre de points fixes de $V$ est ind\'ependant de $\ell$. Autre myst\`ere: pourquoi, le genre $g$ \'etant fix\'e, le nombre de points fixes de $V$ admet-il une description uniforme en terme de $H^1(X, \Ql)$, de sa structure symplectique \`a valeurs dans $\Ql(-1)$ et de son endomorphisme $\Frob^*$, alors que $\pi_1(X,x)$ est un quotient du compl\'et\'e profini du $\pi_1$ topologique de $\mathcal{P}$, quel quotient d\'ependant de $X$.

\subsection{}
Soit $E_r$ l'ensemble des classes d'isomorphie de $\Ql$-faisceaux lisses irr\'eductibles de rang $r$ sur $X$. Gr\^ace \`a Lafforgue [L], le nombre de points fixes de la bijection $\Frob^*: E_r \to E_r$ a une interpr\'etation automorphe. Pour $S_0 \subset X_0$ un ensemble fini de points ferm\'es de $X_0$, et $S = S_0 \otimes_{\F_q} F$, on peut aussi consid\'erer des $\Ql$-faisceaux lisses irr\'eductibles sur $X - S$, avec ramification prescrite en chaque $s \in S$. Dans tous les cas qui ont \'et\'e calcul\'es, les miracles (A) (B) d\'ecouverts par Drinfeld persistent, mutatis mutandis. Le but de l'expos\'e est de donner un panorama de ce qui est connu.

\section{Que compter?}

\subsection{}
Fixons un entier $r \geq 1$, et soient $X_0, S_0/\F_q$ et $X, S/F$ comme en 1.1 et 1.7. Pour $s \in S$, soient $X_{(s)}$ l'hens\'elis\'e de $X$ en $s$ et $X_{(*s)} := X_{(s)} - \{s\}$. Supposons donn\'e, pour chaque $s \in S$, un $\Ql$-faisceau lisse de rang $r$ $R(s)$ sur $X_{(*s)}$. Seule sa classe d'isomorphie importe. Le morphisme $\Frob: X \to X$ induit un morphisme encore not\'e $\Frob$ de $X_{(*s)}$ dans $X_{(*\Frob(s))}$. On suppose que, pour tout $s$ dans $S$, $R(s)$ est isomorphe \`a l'image inverse de $R(\Frob(s))$ par $\Frob$:
\begin{equation}
R(s) \xrightarrow{\sim} \Frob^*(R(\Frob(s))).
\tag{2.1.1}
\end{equation}

Soit $R := (R(s))_{s \in S}$ la famille des $R(s)$ et notons $E(R)$ l'ensemble des classes d'isomorphie de $\Ql$-faisceaux lisses irr\'eductibles de rang $r$ $F$ sur $X - S$, tels que pour tout $s \in S$, l'image inverse $F|_{X_{(*s)}}$ de $F$ sur $X_{(*s)}$ soit isomorphe \`a $R(s)$. L'hypoth\`ese (2.1.1) assure que le foncteur image inverse par $\Frob$ induit une permutation de $E(R)$. Nous noterons $V$ cette permutation.

Notre choix de ``$V$'', premi\`ere lettre de ``Verschiebung'' est d\^u \`a une analogie avec le cas de certaines vari\'et\'es ab\'eliennes, expliqu\'ee en 2.2 qui suit. Dans cette analogie, $q$ est classiquement un nombre premier, plut\^ot qu'une puissance d'un nombre premier.

\subsection{}
Pour tout sch\'ema $Y$ sur $\F_q$, notons $\Frob_Y$ l'endomorphisme du $\F_q$-sch\'ema $Y$ qui est l'identit\'e sur l'espace topologique sous-jacent, et qui est donn\'e sur le faisceau structural par $\Frob_Y^*(f) = f^q$. C'est l'endomorphisme de Frobenius de $Y$.

Prenons pour $Y$ la vari\'et\'e ab\'elienne $\Pic^0(X_0)$ sur $\F_q$. On obtient
\begin{equation}
\Frob_{\Pic^0(X_0)}: \Pic^0(X_0) \longrightarrow \Pic^0(X_0).
\tag{2.2.1}
\end{equation}

La jacobienne $\Pic^0(X_0)$ repr\'esente le foncteur suivant sur les $\F_q$-sch\'emas: le faisceau (pour la topologie de Zariski, ou pour la topologie \'etale, cela revient au m\^eme) associ\'e au pr\'efaisceau qui \`a $S$ attache le groupe $\Pic^0(S \times_{\F_q} X_0)$ des faisceaux inversibles sur $S \times_{\F_q} X_0$, de degr\'e z\'ero sur chaque fibre de la projection sur $S$. Si $f: Y_0 \to X_0$ est un morphisme, l'image inverse de faisceaux inversibles par $f$ d\'efinit un morphisme $f^*: \Pic^0(X_0) \to \Pic^0(Y_0)$. Prenons pour $f$ l'endomorphisme $\Frob_{X_0}$ de $X_0$. L'endomorphisme $\Frob_{X_0}^*$ de $\Pic^0(X_0)$ m\'erite d'\^etre appel\'e $V$, car il v\'erifie
\begin{equation}
\Frob_{X_0}^* \circ \Frob_{\Pic^0(X_0)} = \Frob_{\Pic^0(X_0)} \circ \Frob_{X_0}^* = q.
\tag{2.2.2}
\end{equation}

Preuve de (2.2.2). Interpr\'etons (2.2.2) comme deux \'egalit\'es entre endomorphismes du foncteur repr\'esent\'e par $\Pic^0(X_0)$. Pour tout foncteur contravariant $G$ sur la cat\'egorie des $\F_q$-sch\'emas, l'endomorphisme de Frobenius $\Frob_G$ de $G$ est d\'efini par $G \xrightarrow{\Frob_G(S): G(S) \to G(S),\text{ est } G(\Frob_S)}$. Si $G$ repr\'esente $Y$, $\Frob_G$ est induit par $\Frob_Y$.

Prenons pour $G$ le foncteur $S \mapsto \Pic(S \times_{\F_q} X_0)$ des classes d'isomorphie de faisceaux inversibles sur $S \times_{\F_q} X_0$. L'endomorphisme $\Frob_G(S)$ de $\Pic(S \times_{\F_q} X_0)$ est l'image inverse de classes d'isomorphie de faisceaux inversibles par $\Frob_S \times \Id: S \times_{\F_q} X_0 \to S \times_{\F_q} X_0$.

Si on compose cette image inverse avec l'image inverse par $\Frob_{X_0}$, on obtient l'image inverse par $\Frob_{S \times_{\F_q} X_0}$, qui n'est autre que l'\'el\'evation \`a la puissance tensorielle $q$. Ne consid\'erant que les faisceaux inversibles de degr\'e z\'ero sur les fibres de la projection sur $S$, et passant du pr\'efaisceau obtenu au faisceau associ\'e, on obtient (2.2.2).

Si maintenant on \'etend les scalaires de $\F_q$ \`a $F$, $\Frob_{X_0}$ devient l'endomorphisme $\Frob$ de $X$, $\Pic^0(X_0)$ devient $\Pic^0(X)$ et l'endomorphisme $\Frob_{X_0}^*$ de $\Pic^0(X_0)$ devient l'endomorphisme de $\Pic^0(X)$ ``image inverse par $\Frob$''.

\subsection{}
Les questions que nous nous posons sont les suivantes.

\begin{enumerate}[label=(\roman*)]
\item Calculer, pour $n \geq 1$, le nombre $N_n(R)$ de points fixes de l'it\'er\'e $n$i\`eme $V^n$ de la permutation $V$ de $E(R)$ (2.1). On notera que remplacer $(X_0, S_0)$ par $(X_0, S_0) \otimes_{\F_q} \F_{q^n}$ remplace $V$ par $V^n$.
\item D\'eterminer si les miracles 1.2 (A) (B) subsistent.
\item Si oui, pourquoi?
\end{enumerate}

\subsection{}
Soit $F$ un $\Ql$-faisceau lisse sur $X - S$. Comme en 1.2, on a

\begin{lemme}
\begin{enumerate}[label=(\roman*)]
\item Pour que la classe d'isomorphie de $F$ soit fixe sous $\Frob^*$, il faut et il suffit que $F$ soit isomorphe \`a l'image inverse sur $X$ d'un $\Ql$-faisceau $F_0$ sur $X_0$.
\item De plus, si $F$ est irr\'eductible, un tel $F_0$ est unique \`a $\F_q$-torsion pr\`es, et
\item les $\F_q$-tordus de $F_0$ par $\lambda \in \Zl^*$ (1.2) sont deux \`a deux non isomorphes.
\end{enumerate}
\end{lemme}

D'apr\`es 2.4, le nombre de points fixes de la permutation $V$ de $E(R)$ est le nombre de classes, modulo $\F_q$-torsion, de $\Ql$-faisceaux lisses de rang $r$ $F_0$ sur $X_0 - S_0$, tels que
\begin{enumerate}[label=(\alph*)]
\item l'image inverse $F$ de $F_0$ sur $X - S$ est irr\'eductible;
\item pour chaque $s \in S$, $F|_{X_{(*s)}}$ est isomorphe \`a $R(s)$.
\end{enumerate}

Soit $D_0 = \sum n_i x_i$ un diviseur de degr\'e un de $X_0 - S_0$. Regardons $F_0$ comme une repr\'esentation de Galois (voir 2.7), soit $\Fr_{x_i}$ un Frobenius en $x_i$, et posons
\[
\langle D_0, F_0 \rangle_i := \det(\Fr_{x_i}, F_0)^{n_i}.
\]
Si $F_0'$ est le $\F_q$-tordu de $F_0$ par $\lambda$, on a
\[
\langle D_0, F_0' \rangle = \lambda^{r \langle D_0, F_0 \rangle}.
\]
D'apr\`es 2.4 (ii) et (iii), il revient donc au m\^eme de compter les $F_0$ v\'erifiant (a) et (b) \`a $\F_q$-torsion pr\`es, ou de compter, \`a isomorphisme pr\`es, seulement ceux qui v\'erifient en outre $\langle D_0, F_0 \rangle = 1$, et de diviser par $r$. Pour une contrepartie automorphe, voir la remarque qui suit 2.6.

\subsection{}
Posons $K := k(X_0)$. Pour chaque point ferm\'e $x$ de $X_0$, soient $k(x)$ le corps r\'esiduel de $x$, $\deg(x) := [k(x): \F_q]$ son degr\'e, $K_x$ le compl\'et\'e de $K$ en $x$, et $\mathcal{O}_x$ l'anneau de la valuation $v_x$ de $K_x$. Soit $\mathbb{A}$ l'anneau des ad\`eles de $K$. La somme $v := \sum \deg(x) v_x: \mathbb{A}^* \to \Z$ est triviale sur $K^*$. On a $|\mathbb{A}^*/K^*|_k = q^{-v(x)}$.

D'apr\`es Lafforgue, $F_0$ correspond \`a une repr\'esentation automorphe cuspidale $\pi$ de $GL(r, \mathbb{A})$. Si $F_0'$ est le $\F_q$-tordu de $F_0$ par $\lambda$, la repr\'esentation $\pi'$ correspondante est la $\F_q$-tordue de $\pi$ par $\lambda$, d\'efinie comme suit. Comme repr\'esentation, c'est le produit tensoriel de $\pi$ par la repr\'esentation de dimension un $g \mapsto \lambda^{v \det g}$. Comme espace de fonctions sur $GL(r, \mathbb{A})$, c'est l'espace des produits de $f$ dans $\pi$ par le caract\`ere $g \mapsto \lambda^{v \det g}$.

Traduit dans le langage automorphe, le probl\`eme 2.3 (i) est celui de compter \`a $\F_q$-torsion pr\`es les repr\'esentations automorphes cuspidales $\pi$ de $GL(r, \mathbb{A})$, non ramifi\'ees en dehors de $S_0$, et telles que
\begin{enumerate}[label=(\alph*$'$)]
\item Pour tout $n \geq 1$, $\pi$ reste cuspidale apr\`es changement de base de $K$ \`a $K \otimes \F_{q^n}$.
\end{enumerate}
Cette condition exprime l'irr\'eductibilit\'e 2.4 (a) de $F$. Il suffit de la v\'erifier pour $n$ divisant $r$.
\begin{enumerate}[label=(\alph*$'$),start=2]
\item Pour chaque $s \in S_0$ et $s' \in S$ au-dessus de $s_0$, la composante locale $\pi_{s_0}$ de $\pi$ v\'erifie une condition dict\'ee par $R(s)$ et par la correspondance de Langlands locale.
\end{enumerate}

\begin{lemme}
Si $\pi$ automorphe cuspidale v\'erifie (a$'$), les $\F_q$-tordues de $\pi$ sont deux \`a deux distinctes.
\end{lemme}
C'est la traduction de 2.4 (iii).

Si $d$ est un id\`ele v\'erifiant $v(d) = 1$, il revient donc au m\^eme de compter \`a $\F_q$-torsion pr\`es les $\pi$ v\'erifiant (a$'$) (b$'$), ou de ne compter que ceux dont le caract\`ere central $\omega_\pi$ v\'erifie $\omega_\pi(d) = 1$, et de diviser par $r$.

En caract\'eristique finie, la notion de repr\'esentation automorphe est purement alg\'ebrique. Ceci permet de prendre les repr\'esentations automorphes $\pi$ comme \'etant \`a coefficients dans $\Ql$, plut\^ot que $\C$. Si on le fait, $\pi_{s_0}$ est une $\Ql$-repr\'esentation admissible irr\'eductible de $GL(r, K_{s_0})$, et la correspondance de Langlands locale lui associe une classe d'isomorphie de $\Ql$-repr\'esentations $F$-semi-simples continues de dimension $r$ du groupe de Weil local $W(\overline{K}_{s_0}/K_{s_0})$. La condition (b$'$) est que sa restriction au sous-groupe d'inertie soit donn\'ee par $R(s)$ (cf 2.4 ci-dessous).

Presque tous les r\'esultats connus sur la question 2.3 (i) sont obtenus via le dictionnaire 2.5.

\subsection{}
Dans 2.1, nous avons employ\'e le langage des $\Ql$-faisceaux. Ceci nous a dispens\'e d'avoir \`a choisir des points base. Un langage \'equivalent est celui des repr\'esentations de Galois. Voici la traduction.

Soit $\overline{k(X)}$ une cl\^oture alg\'ebrique du corps des fonctions rationnelles $k(X)$ de $X$, et $k(X)^{nr}$ la plus grande sous-extension non ramifi\'ee en dehors de $S$. Le groupe de Galois $\Gal(k(X)^{nr}/k(X))$ est le groupe fondamental de $X$ en le point g\'eom\'etrique $\overline{\eta} := \Spec(\overline{k(X)})$ de $X$: le foncteur ``fibre en $\overline{\eta}$'' est une \'equivalence de cat\'egories, des $\Ql$-faisceaux lisses (resp.\ et irr\'eductibles de rang $r$) sur $X - S$, vers les $\Ql$-repr\'esentations lin\'eaires de $\Gal(k(X)^{nr}/k(X))$ (resp.\ irr\'eductibles de rang $r$).

Le sch\'ema $X_{(*s)}$ est le spectre d'une extension $k(X_{(*s)})$ de $k(X)$. Si on choisit un plongement de cette extension dans $\overline{k(X)}$, le groupe d'inertie correspondant $I_s := \Gal(\overline{k(X)}/k(X_{(*s)}))$ est de m\^eme le groupe fondamental de $X_{(*s)}$ en $\overline{\eta}$. Il s'envoie dans $\pi_1(X, \overline{\eta})$:
\begin{equation}
\pi_1(X_{(*s)}, \overline{\eta}) = I_s = \Gal(\overline{k(X)}/k(X_{(*s)})) \hookrightarrow \Gal(k(X)^{nr}/k(X)) = \pi_1(X, \overline{\eta}).
\tag{2.7.1}
\end{equation}
La fibre en $\overline{\eta}$ de $R(s)$ est une repr\'esentation $\rho_s$ de $I_s$.

Soit $F$ un $\Ql$-faisceau lisse sur $X - S$, d'image inverse $F|_{X_{(*s)}}$ sur $X_{(*s)}$. Les fibres au point g\'eom\'etrique $\overline{\eta}$ sont respectivement une repr\'esentation de $\pi_1(X - S, \overline{\eta}) = \Gal(k(X)^{nr}/k(X))$ et sa restriction \`a $\pi_1(X_{(*s)}, \overline{\eta}) = I_s$ par (2.7.1). Que $F|_{X_{(*s)}}$ soit isomorphe \`a $R(s)$ \'equivaut \`a ce que cette restriction \`a $I_s$ soit isomorphe \`a $\rho_s$.

La condition (2.1.1) \'equivaut \`a ce que les $\rho_s$, pour $s$ au-dessus de $s_0$, proviennent d'une repr\'esentation d'un groupe de d\'ecomposition $D_{s_0} \subset \Gal(\overline{k(X)}/k(X_0))$ en $s_0$. Elle implique que les repr\'esentations $\rho_s$ sont quasi-unipotentes.

\subsection{}
J'esp\`ere que le probl\`eme 2.3 (i) est plus maniable lorsqu'un $\Ql$-faisceau lisse $F$ tel que $F|_{X_{(*s)}} \cong R(s)$ pour tout $s \in S$ est automatiquement irr\'eductible. C'est le cas si la condition suivante est v\'erifi\'ee. Consid\'erons un entier $0 < r' < r$ et la donn\'ee, pour chaque $s \in S$, d'un sous-$\Ql$-faisceau $R'(s)$ lisse de rang $r'$, de $R(s)$. Posons $R''(s) := R(s)/R'(s)$. La condition est que
\begin{multline}
\text{Quels que soient } r' \text{ et les } R'(s), \text{ ou bien il n'existe aucun } \Ql\text{-faisceau } F' \text{ lisse de rang } r' \text{ sur } X - S \text{ tel que } F'|_{X_{(*s)}} \cong R'(s) \text{ pour tout } s \in S, \\ \text{ou bien il n'existe aucun } \Ql\text{-faisceau } F'' \text{ lisse de rang } r - r' \text{ sur } X - S \text{ tel que } F''|_{X_{(*s)}} \cong R''(s) \text{ pour tout } s \in S.
\tag{2.8.1}
\end{multline}
La condition (2.8.1) est v\'erifi\'ee si l'un des $R(s)$ est irr\'eductible. Pour d'autres cas, voir 2.10.4, 2.11 et 2.12.

\subsection{}
Pour tout corps $k$ de caract\'eristique $p$ ou $0$, contenant toutes les racines de l'unit\'e d'ordre premier \`a $p$, posons
\begin{equation}
\Z_{p'}(1)(k) := \varprojlim \mu_n(k).
\tag{2.9.1}
\end{equation}
La limite projective est prise selon l'ensemble ordonn\'e par inclusion des sous-groupes $\mu_n(k)$ de $k^*$ ($n$ premier \`a $p$). Le morphisme de transition de $\mu_m(k)$ \`a $\mu_n(k)$ est $x \mapsto x^{m/n}$.

Si cela ne cr\'ee pas d'ambigu\"it\'e, on omettra ``$p'$'' et la mention de $k$.

\subsection{}
\`A cause de (c), on ne peut m\^eme pas esp\'erer que $Z$ provienne d'un champ alg\'ebrique $\mathcal{C}$ (c'est-\`a-dire que $Z(T)$ soit l'ensemble des classes d'isomorphie dans $\mathcal{C}(T)$).

Toutefois
\begin{enumerate}[label=(\alph*)]
\setcounter{enumi}{3}
\item Les groupes d'automorphismes sont des groupes alg\'ebriques connexes. Si $k' \subset k''$ sont des extensions finies de $\F_q$, ceci assure que
\[
Z(k') \xrightarrow{\sim} Z(k'')^{\Gal(k''/k')}.
\]
\item Dans (b) ci-dessus, la fonction $t \mapsto \dim \End(E_t, \alpha_t)$ est constructible; dans (c), l'ensemble des $t$ tels que $(E_t, \alpha_t)$ soit g\'eom\'etriquement ind\'ecomposable est constructible.
\end{enumerate}

Les propri\'et\'es (d) et (e) permettent, en ``d\'ecoupant $Z$ en morceaux'' et en prenant les faisceaux (repr\'esentables) associ\'es, d'obtenir $Z$ tel que promis par (3.3.1).

\subsection{Esquisse de preuve de 3.5 $\Rightarrow$ (3.3.1)}
Pour tout sch\'ema $T$ sur $\F_q$, soient $X_T := X_0 \times_{\F_q} T$ et $S_T := S_0 \times_{\F_q} T$. D\'efinissons $Z(T)$ comme \'etant l'ensemble de classes d'isomorphie de fibr\'es vectoriels de rang $r$ sur $X_T$, munis d'une structure parabolique le long de $S_T$, et tels que les conditions de 3.2 et 3.3 soient v\'erifi\'ees sur chaque fibre g\'eom\'etrique de $X_T \to T$: degr\'e $d$, type de la structure parabolique, ind\'ecomposabilit\'e.

Na\"ivement, on aimerait prendre pour sch\'ema $Z$ de (3.3.1) un sch\'ema qui repr\'esente le foncteur $Z$. Pour plusieurs raisons, ce foncteur n'est pas repr\'esentable:
\begin{enumerate}[label=(\alph*)]
\item Les objets consid\'er\'es ayant des automorphismes non triviaux, par exemple les homoth\'eties, ce foncteur n'est pas un faisceau, m\^eme pour la topologie de Zariski. Pour le foncteur $\Pic$, on contourne le m\^eme probl\`eme en passant au faisceau associ\'e. Ici, la situation est pire, car, pour $(E_T, \alpha_T)$ un fibr\'e vectoriel \`a structure parabolique sur $X_T$, et pour $(E, \alpha)$ le fibr\'e induit sur la fibre $X_t$ de $X_T \to T$ en $t \in T$:
\item la dimension du groupe des automorphismes de $(E_t, \alpha_t)$ (\'egale \`a $\dim \End(E_t, \alpha_t)$) n'est pas toujours localement constante en $t$.
\item L'ensemble des $t$ pour lesquels $(E_t, \alpha_t)$ est g\'eom\'etriquement ind\'ecomposable n'est pas toujours localement ferm\'e.
\end{enumerate}
\`A cause de (c), on ne peut m\^eme pas esp\'erer que $Z$ provienne d'un champ alg\'ebrique $\mathcal{C}$ (c'est-\`a-dire que $Z(T)$ soit l'ensemble des classes d'isomorphie dans $\mathcal{C}(T)$).

Toutefois
\begin{enumerate}[label=(\alph*)]
\setcounter{enumi}{3}
\item Les groupes d'automorphismes sont des groupes alg\'ebriques connexes. Si $k' \subset k''$ sont des extensions finies de $\F_q$, ceci assure que
\[
Z(k') \xrightarrow{\sim} Z(k'')^{\Gal(k''/k')}.
\]
\item Dans (b) ci-dessus, la fonction $t \mapsto \dim \End(E_t, \alpha_t)$ est constructible; dans (c), l'ensemble des $t$ tels que $(E_t, \alpha_t)$ soit g\'eom\'etriquement ind\'ecomposable est constructible.
\end{enumerate}

Les propri\'et\'es (d) et (e) permettent, en ``d\'ecoupant $Z$ en morceaux'' et en prenant les faisceaux (repr\'esentables) associ\'es, d'obtenir $Z$ tel que promis par (3.3.1).

\subsection{Esquisse de preuve de 3.5 (pour $n = 1$)}
On applique la strat\'egie 2.5. Il nous faut montrer que $|Z(\F_q)|$ est le nombre de repr\'esentations automorphes cuspidales $\pi$ de $GL(r, \mathbb{A})$, compt\'ees \`a $\F_q$-torsion pr\`es, qui sont non ramifi\'ees en dehors de $S_0$ et dont la composante locale $\pi_{s_0}$ en $s \in S_0$ est du type suivant. Par la correspondance de Langlands locale, la repr\'esentation $\pi_{s_0}$ de $GL(r, K_{s_0})$ correspond \`a une $\Ql$-repr\'esentation $F$-semi-simple $\rho$ du groupe de Weil $W(\overline{K}_{s_0}/K_{s_0})$. On veut en premier lieu que $\rho$ soit semi-simple (dans le langage du groupe de Weil-Deligne: $N = 0$; dans le langage d'Arthur: action triviale du $SL(2)$ ou $SU(2)$ local). Le groupe $W(\overline{K}_{s_0}/K_{s_0})^{ab}$ est $K_{s_0}^*$. Son groupe d'inertie $O_{s_0}^*$ admet pour quotient $k(s_0)^*$. On veut que la restriction de $\rho$ au groupe d'inertie se factorise par la repr\'esentation de $k(s_0)^*$ somme des $\varepsilon$ dans $R(s_0)$.

Soit $P[s_0] \subset GL(r, \mathcal{O}_{s_0})$ le sous-groupe parahorique image inverse du parabolique standard $\overline{P}[s_0]$ de $GL(r, k(s_0))$ de type $n(s_0)$. Le quotient r\'eductif $L[s_0]$ de $\overline{P}[s_0]$ est $\prod_i GL(n[s_0](i), k(s_0))$. Notons $U[s_0]$ le noyau de $P[s_0] \to \overline{P}[s_0] \to L[s_0]$ et $\varepsilon[s_0]$ le caract\`ere
\[
\varepsilon[s_0] = \prod_i \varepsilon[s_0](i)(\det g_i): P[s_0] \twoheadrightarrow \overline{P}[s_0] \twoheadrightarrow L[s_0] \twoheadrightarrow \Ql^*
\]
du parahorique $P[s_0]$. Par construction, les $\varepsilon[s_0](i)$ sont deux \`a deux distincts.

Nous admettrons l'assertion suivante:

\textbf{Assertion 3.8.} Les repr\'esentations admissibles irr\'eductibles de $GL(r, K_{s_0})$ du type d\'efini en 3.7 sont celles qui admettent une droite stable par $P$ sur laquelle $P$ agit par le caract\`ere $\varepsilon[s_0]$. Cette droite est unique.

\subsection{}
Soit $\mathcal{A}'$ l'ensemble des repr\'esentations automorphes cuspidales $\pi$ de $GL(r, \mathbb{A})$, admettant une droite $L(\pi)$ fixe par les $GL(r, \mathcal{O}_x)$ pour $x \notin S_0$, et stable par $P[s_0]$ pour $s_0 \in S_0$, l'action de $P[s_0]$ \'etant donn\'ee par $\varepsilon[s_0]$. La droite $L(\pi)$ est unique si elle existe.

D'apr\`es Lafforgue [L] et 3.8, $N_1(R)$ est le nombre de classes modulo $\F_q$-torsion dans $\mathcal{A}'$. Fixons un id\`ele $d$ de valuation $1$ et soit $\mathcal{A}$ l'ensemble des $\pi$ dans $\mathcal{A}'$ dont le caract\`ere central $\omega_\pi$ vaut $1$ en $d$. Comme observ\'e apr\`es 2.3.1, on a
\begin{equation}
N_1(R) = |\mathcal{A}' \text{ mod } \F_q\text{-torsion}| = |\mathcal{A} \text{ mod } \F_q\text{-torsion}| \cdot \frac{1}{r} = \frac{|\mathcal{A}|}{r}.
\tag{3.9.1}
\end{equation}

\subsection{}
Soit $\mathcal{A}$ l'espace vectoriel des fonctions sur $GL(r, \mathbb{A})$ qui sont
\begin{enumerate}[label=(\alph*)]
\item invariantes \`a gauche par $GL(r, K) \cdot d^\Z$;
\item invariantes \`a droite sous $GL(r, \mathcal{O}_x)$, pour $x \notin S_0$;
\item se transforment par $\varepsilon[s_0]$ pour l'action \`a droite de $P[s_0]$, pour $s_0 \in S_0$.
\end{enumerate}

Un point cl\'e de la preuve de 3.5 est le lemme suivant.

\begin{lemme}
L'espace vectoriel $\mathcal{A}$ est la somme directe des droites $L(\pi)$, pour $\pi$ dans $\mathcal{A}$. Il est donc de dimension $r N_1(R)$.
\end{lemme}

Ce lemme est la contrepartie automorphe, et une cons\'equence, de ce que l'irr\'eductibilit\'e des $\Ql$-faisceaux $F_0$ consid\'er\'es r\'esulte des hypoth\`eses locales faites.

Il s'agit de v\'erifier que les fonctions dans $\mathcal{A}$ sont automatiquement cuspidales. Ceci fait, si on \'ecrit l'espace de fonctions cuspidales sur $GL(r, K) \cdot d^\Z \backslash GL(r, \mathbb{A})$ comme somme directe de repr\'esentations irr\'eductibles $\pi$ de $GL(r, \mathbb{A})$, les $\pi$ dans $\mathcal{A}$ contribuent $L(\pi)$ \`a $\mathcal{A}$, les autres ne contribuent pas.

Que les fonctions dans $\mathcal{A}$ soient cuspidales r\'esulte de ce que dans la d\'ecomposition spectrale des fonctions sur $GL(r, K) \cdot d^\Z \backslash GL(r, \mathbb{A})$, les repr\'esentations induites donnant lieu aux s\'eries d'Eisenstein n'admettent pas de vecteur non nul v\'erifiant 3.10 (b) (c).

\subsection{}
L'espace $\mathcal{A}$ somme des $L(\pi)$ pour $\pi \in \mathcal{A}$ s'identifie \`a l'espace des fonctions $f$ sur
\begin{equation}
GL(r, K) d^\Z \backslash GL(r, \mathbb{A}) / \prod_{x \notin S_0} GL(r, \mathcal{O}_x) \prod_{s_0 \in S_0} U[s_0]
\tag{3.12.1}
\end{equation}
telles que, pour chaque $s \in S_0$ l'action \`a droite de $P[s_0]$ v\'erifie $f(gp) = \varepsilon[s_0](p) f(g)$.

Notons $\overline{\mathcal{A}}$ le quotient de $\mathcal{A}$ par l'action par $\F_q$-torsion de $\mu_r(\Ql)$, et pour chaque orbite $B$, soit $V_B$ le sous-espace de $\mathcal{A}$ somme des $L(\pi)$ pour $\pi$ dans $B$. D'apr\`es (2.3.1), on a $\dim(V_B) = r$.

Pour $m \in \Z/r$, notons $(3.12.1)_m$ la partie de (3.12.1) o\`u $v \det = m$, $\varphi_m$ sa fonction caract\'eristique et $L(\pi, m) = \varphi_m \cdot f$ pour $f \in L(\pi)$. La dimension de $L(\pi, m)$ est $\leq 1$. Pour chaque caract\`ere $\eta$ de $\Z/r$, si $\pi'$ est le $\F_q$-tordu de $\pi$ par $\eta(1)$, la multiplication par $\eta^v$ transforme $L(\pi)$ en $L(\pi')$. Pour $\pi \in B$, $L(\pi, m)$ est donc ind\'ependant de $\pi$; on pose $L(B, m) := L(\pi, m)$. L'espace $V_B$ est stable par les multiplications par les $\eta^v$, donc par multiplication par les $\varphi_m$: pour $\pi \in B$,
\[
V_B = \bigoplus_{m \in \Z/r} L(B, m).
\]
Chaque $L(B, m)$ est donc de dimension un.

La somme $\overline{\mathcal{A}}$ des $L(B, m)$ est de dimension $|\mathcal{A}|/r = N_1(R)$. Elle s'identifie \`a l'espace des fonctions $f_m$ sur (3.12.1) telles que, pour chaque $s \in S_0$, l'action \`a droite de $P[s_0]$ v\'erifie $f(gp) = \varepsilon[s_0](p) f(g)$.

\begin{lemme}
$N_1(R)$ est la dimension de l'espace, encore not\'e $\mathcal{A}_d$, des fonctions $f$ sur
\begin{equation}
GL(r, K) \backslash GL(r, \mathbb{A})(d) / \prod_{x \notin S_0} GL(r, \mathcal{O}_x) \prod_{s_0 \in S_0} U[s_0]
\tag{3.12.2}
\end{equation}
qui, pour l'action \`a droite des $P[s_0]$ ($s \in S_0$) v\'erifient $f(gp) = \varepsilon[s_0](p) f(g)$.
\end{lemme}

\subsection{}
L'ensemble (3.12.2) s'interpr\`ete comme l'ensemble des classes d'isomorphie de triples $(E_0, \alpha, \beta)$ o\`u $E_0$ est un fibr\'e vectoriel de rang $r$ et de degr\'e $d$ sur $X_0$, et o\`u $\alpha$ et $\beta$ sont des structures des types suivants sur $E_0$:
\begin{enumerate}[label=(\alph*)]
\item $\alpha$ est une structure parabolique de type $(n[s_0])_{s_0 \in S_0}$;
\item $\beta$ est la donn\'ee de bases des espaces vectoriels $\Gri_{F(s_0)}(E_{s_0})$.
\end{enumerate}
Nous identifierons une base comme en (b) avec un isomorphisme
\[
\beta: \bigoplus_{s_0, i} k(s_0)^{n[s_0](i)} \xrightarrow{\sim} \Gri_{F(s_0)}(E_{s_0}).
\]
Avec cette interpr\'etation, l'espace $\mathcal{A}_d$ de 3.13 devient l'espace des fonctions $f(E_0, \alpha, \beta)$, pour $E_0, \alpha, \beta$ comme ci-dessus, telles que
\begin{enumerate}[label=(\roman*)]
\item $f(E_0, \alpha, \beta)$ ne d\'epend que de la classe d'isomorphisme de $(E_0, \alpha, \beta)$;
\item pour $s \in S_0$ et $g \in L[s_0] = \prod_i GL(n, k(s_0))$, on a
\[
f(E_0, \alpha, \beta g) = \varepsilon[s_0](g) f(E_0, \alpha, \beta).
\]
\end{enumerate}

Pour $T$ un automorphisme de $(E_0, \alpha)$, notons $\Gri_{F(s_0)}(T_{s_0})$ l'automorphisme de $\Gri_{F(s_0)}(E_{s_0})$ induit par $T$, et posons
\begin{equation}
\chi(T) := \prod_{s_0} \prod_i \varepsilon[s_0](i)(\det \Gri_{F(s_0)}(T_{s_0})).
\tag{3.14.1}
\end{equation}

Quel que soit $\beta$, l'hypoth\`ese (i) donne que $f(E_0, \alpha, T(\beta)) = f(E_0, \alpha, \beta)$. L'hypoth\`ese (ii) donne que $f(E_0, \alpha, T(\beta)) = \chi(T) f(E_0, \alpha, \beta)$. S'il existe un automorphisme $T$ tel que $\chi(T) \neq 1$, on a donc $f(E_0, \alpha, \beta) = 0$. Par contre, si pour tout automorphisme $T$ de $(E_0, \alpha)$ on a $\chi(T) = 1$, $f$ peut \^etre librement prescrit en un $(E_0, \alpha, \beta)$, et la valeur en $(E_0, \alpha, \beta)$ d\'etermine la valeur en les $(E_0', \alpha, \beta')$.

La dimension de $\mathcal{A}_d$ est donc le nombre de classes d'isomorphie de $(E_0, \alpha)$ comme ci-dessus, tels que pour tout automorphisme $T$ on ait $\chi(T) = 1$. Le th\'eor\`eme 3.5 r\'esulte d\`es lors de la caract\'erisation (3.4.3) de l'ind\'ecomposabilit\'e g\'eom\'etrique et de la proposition suivante.

\begin{proposition}
Soit $E_0$ un fibr\'e vectoriel de rang $r$ sur $X_0$, muni d'une structure parabolique $\alpha$ de type $(n(s_0))_{s_0 \in S_0}$. Les conditions suivantes sont \'equivalentes:
\begin{enumerate}[label=(\roman*)]
\item Le quotient de la $\F_q$-alg\`ebre $\End(E_0, \alpha)$ par son radical est r\'eduit \`a $\F_q$.
\item Quel que soit l'automorphisme $T$ de $(E_0, \alpha)$, $\chi(T)$ d\'efini par (3.14.1) vaut $1$.
\end{enumerate}
\end{proposition}

Nous d\'eduirons que (i) implique (ii) de (3.2.2), et que la n\'egation de (i) implique la n\'egation de (ii) de l'hypoth\`ese de position g\'en\'erale (2.10.3).

Preuve de (i) $\Rightarrow$ (ii). Si $\lambda$ est l'image dans $\F_q^*$ de l'automorphisme $T \in \End(E_0, \alpha)^*$, on a
\[
\prod_i \varepsilon[s_0](i)(\det \Gri_{F(s_0)}(T_{s_0})) = \prod_i \varepsilon(s_0)(i)^{n[s_0](i)(\lambda)},
\]
de sorte que (3.2.2) implique que $\chi(T) = 1$.

Preuve de (ii) $\Rightarrow$ (i). Le quotient de la $\F_q$-alg\`ebre $\End(E_0, \alpha)$ par son radical est un produit d'alg\`ebres de matrices sur des extensions finies de $\F_q$. S'il n'est pas r\'eduit \`a $\F_q$, il contient une $\F_q$-alg\`ebre commutative s\'eparable de dimension $> 1$. Cette alg\`ebre admet un rel\`evement dans $\End(E_0, \alpha)$.

Supposons donc que $\End(E_0, \alpha)$ contienne une sous-alg\`ebre commutative s\'eparable $k$ de dimension $> 1$. Il nous faut montrer que le caract\`ere (3.14.1) de $\End(E_0, \alpha)^*$ est non trivial. Nous prouverons que sa restriction $\chi$ \`a $k^*$ est non triviale.

\textbf{Cas d\'eploy\'e.} Supposons tout d'abord que $k$ est un produit de copies de $\F_q$: $k \cong \F_q^I$, et que les points ferm\'es $s \in S_0$ sont de degr\'e $1$: $\F_q \xrightarrow{\sim} k(s_0)$. La structure de $\F_q^I$-module de $E_0$ fournit une d\'ecomposition
\begin{equation}
E_0 = \bigoplus_{\iota \in I} E_0^\iota
\tag{3.15.1}
\end{equation}
telle que $(\lambda_\iota)_{\iota \in I} \in \F_q^I$ agisse sur $E_0^\iota$ par multiplication par $\lambda_\iota$. La d\'ecomposition (3.15.1) est compatible \`a la structure parabolique. Pour $\iota \in I$, posons
\[
n_\iota[s_0](i) = \dim \Gri_{F(s_0)}((E_0^\iota)_{s_0}).
\]
La somme sur $\iota$ des $n_\iota[s_0](i)$ est le rang du fibr\'e $E_0^\iota$. La restriction de $\chi$ au facteur d'indice $\iota$ de $(\F_q^I)^* = \prod_{\iota \in I} \F_q^*$ est
\begin{equation}
\lambda \mapsto \prod_{s_0 \in S_0} \prod_i \varepsilon[s_0](i)^{n_\iota[s_0](i)(\lambda)}.
\tag{3.15.2}
\end{equation}
Son compos\'e avec le morphisme $\Z(1) \twoheadrightarrow \F_q^*$ de 3.1 est le caract\`ere
\[
\prod_{s \in S} \prod_i \varepsilon[s](i)^{n_\iota[s_0](i)}
\]
de $\Z(1)$. L'hypoth\`ese (2.10.3) assure que ce caract\`ere est non trivial.

\textbf{R\'eduction du cas g\'en\'eral au cas d\'eploy\'e.} Il suffit de montrer que la non-trivialit\'e de $\chi$ est invariante par une extension des scalaires de $\F_q$ \`a une extension finie $\F_{q^n}$. Par une telle extension des scalaires, $X_0$ devient une courbe $X'$ sur $\F_{q^n}$, $S_0$ devient un diviseur $S'$ de $X'$, et $k$ devient $k' := k \otimes_{\F_q} \F_{q^n}$. Pour $s \in S'$ au-dessus de $s_0 \in S_0$, posons $n[s'] := n[s_0]$ et notons $\varepsilon[s](i)$ le compos\'e de $\varepsilon[s'](i)$ avec la norme $N_{k(s')/k(s_0)}$. Le fibr\'e $E'$ sur $X'$ image inverse de $E_0$ h\'erite d'une structure parabolique $\alpha'$ de type $(n(s'))_{s' \in S'}$. Si on r\'ep\`ete pour $E'$ sur $X'/\F_{q^n}$ la construction (3.14.1), on obtient un caract\`ere $\chi'$ de $k'^*$. La courbe $X/F$, les entiers $n[s](i)$ et les caract\`eres $\varepsilon[s](i)$ de $\Z(1)$ proviennent aussi bien de $X_0/\F_q$ que de $X'/\F_{q^n}$.

Pour prouver que la trivialit\'e de $\chi$ \'equivaut \`a celle de $\chi'$, il suffit de prouver que
\begin{equation}
\chi' = \chi \circ N_{k'/k}.
\tag{3.15.3}
\end{equation}
La norme $N_{k'/k}$ est en effet surjective. L'identit\'e (3.15.3) r\'esulte de l'identit\'e (3.16.1) qui suit, appliqu\'ee \`a chaque $\Gri_{F(s_0)}(E_{s_0})$, muni de sa structure de $(k, k(s_0))$-bimodule.

Soient $k_0$ une extension finie de $\F_q$, $N$ un $k_0$-espace vectoriel de dimension finie, $k$ un produit d'extensions finies de $\F_q$ et $\rho: k \to \End_{k_0}(N)$. Le $(k, k_0)$-bimodule $N$ d\'efinit un homomorphisme
\[
[N]: k^* \longrightarrow k_0^*: \quad \lambda \mapsto \det(\rho(\lambda)).
\]
Etendons les scalaires de $\F_q$ \`a $\F_{q^n}$. On obtient $k_0'$, $k'$ et un $(k', k_0')$-bimodule $N'$. Ce bimodule d\'efinit
\[
[N']: k'^* \longrightarrow k_0'^*: \quad \lambda \mapsto \det_{k_0'}(\rho(\lambda)).
\]

\begin{lemme}
Le diagramme
\begin{equation}
\begin{tikzcd}
k'^* \arrow[r, "{[N']}"] \arrow[d, "N_{k'/k}"'] & k_0'^* \arrow[d, "N_{k_0'/k_0}"] \\
k^* \arrow[r, "{[N]}"'] & k_0^*
\end{tikzcd}
\tag{3.16.1}
\end{equation}
est commutatif.
\end{lemme}

Preuve. Etendons les scalaires de $\F_q$ \`a $F$. Le diagramme (3.16.1) se plonge dans un diagramme (3.16.2) analogue, o\`u $k$ et $k_0$ sont remplac\'es par $k_F := k \otimes_{\F_q} F$, $k_{0F} := k_0 \otimes_{\F_q} F$, o\`u $N'$ est remplac\'e par un $(k_F, k_{0F})$-bimodule $N$, et o\`u l'extension $\F_{q^n}$ de $\F_q$ est remplac\'ee par un produit $F^J$ de $n$ copies de $F$.
\[
\begin{tikzcd}
(k_F \otimes F^J)^* \arrow[r, "{[N_F]}"] \arrow[d] & (k_{0F} \otimes F^J)^* \arrow[d] \\
k_F^* \arrow[r, "{[N]}"'] & k_{0F}^*
\end{tikzcd}
\tag{3.16.2}
\]
On a $(k_F \otimes F^J)^* = (k_F^*)^J$, $(k_{0F} \otimes F^J)^* \cong (k_{0F}^*)^J$ et les applications verticales ``norme'' sont les applications produit. La premi\`ere ligne horizontale est le produit de copies index\'ees par $J$ de la seconde ligne horizontale. La commutativit\'e de 3.16.2 en r\'esulte.

\section{Formule des traces}

\subsection{}
Dans [DF], le probl\`eme de comptage 2.3 (i) est r\'esolu lorsque $|S_0| \geq 2$ et que les monodromies locales impos\'ees sont unipotentes, avec un seul bloc de Jordan. Dans le langage automorphe, cela signifie qu'on demande aux composantes locales $\pi_{s_0}$ ($s_0 \in S_0$) d'\^etre de la forme
\begin{equation}
\text{repr\'esentation sp\'eciale } \otimes \chi \det g,
\tag{4.1.1}
\end{equation}
pour $\chi$ un caract\`ere non ramifi\'e de $K_{s_0}^*$.

L'hypoth\`ese $|S_0| \geq 2$ permet de passer des repr\'esentations automorphes pour $GL(r, K)$ aux repr\'esentations automorphes pour le groupe multiplicatif d'une alg\`ebre \`a division $D$ de dimension $r^2$ sur $K$. Il est rassurant de n'utiliser la formule des traces que dans un cas \`a quotient compact (modulo le centre), mais il devrait \^etre possible d'utiliser plut\^ot la formule des traces pour $GL(n)$, pour une fonction test convenable, avec la simplification qu'apporte le fait qu'on veut d\'etecter des repr\'esentations automorphes qui en deux places sont de la s\'erie discr\`ete. Apr\`es tout, c'est ainsi qu'on relie $GL(r, K)$ et $D^*$. On simplifiera les explications qui suivent en restant avec $GL(r, K)$, mais en admettant que les seuls termes non nuls dans la formule des traces requise sont ceux associ\'es aux classes de conjugaison d'\'el\'ements elliptiques d'ordre fini de $GL(r, K)$.

Choisissons un plongement $\F_{q^r} \hookrightarrow M_r(\F_q)$. Il d\'efinit un morphisme $\F_{q^r}^* \hookrightarrow GL(r, \F_q) \hookrightarrow GL(r, K)$, par lequel l'ensemble des orbites de $\Gal(\F_{q^r}/\F_q)$ dans $\F_{q^r}^*$ s'envoie bijectivement sur l'ensemble des classes de conjugaison \`a consid\'erer.

Soit $T(\gamma)$ le terme de la formule des traces (donnant le nombre de points fixes de $V$) ainsi associ\'e \`a $\gamma$ dans $\F_{q^r}^*$. Il ne d\'epend que de la sous-extension $\F_{q^m}$ de $\F_q$ engendr\'ee par $\gamma$. Posons $T(m) := T(\gamma)$. Le nombre $c_m$ d'\'el\'ements de $\F_{q^m}^*$ qui engendrent l'extension $\F_{q^m}$ de $\F_q$ est
\[
c_m = \sum_{a|m} \mu(a)(q^{m/a} - 1).
\]

Le nombre de points fixes de $V$ est
\begin{equation}
N_1 = \sum_{m|r} c_m T(m) = \sum_{m|r} \sum_{ab=m} \mu(a)(q^b - 1) T(ab)
\tag{4.1.2}
\end{equation}
\[
\phantom{N_1} = \sum_{b|ab|r} \sum_{a} \mu(a)(q^b - 1) T(ab).
\]

Le nombre de points fixes de $V^n$ ($n \geq 1$) est obtenu de m\^eme, apr\`es avoir \'etendu les scalaires de $\F_q$ \`a $\F_{q^n}$.

\begin{surprise}
Quand on laisse varier $n$, chaque $T(m)$ n'est en g\'en\'eral pas, comme fonction de $n$, de la forme (1.2.1). Par contre, dans (4.1.2), chaque terme de la somme sur $b$ est de cette forme.
\end{surprise}

Je n'ai aucune explication, ni g\'eom\'etrique, ni automorphe, de 4.2, seulement une d\'emonstration. Je n'ai pas non plus d'explication \`a la

\begin{surprise}
Sous les hypoth\`eses de 4.1, le nombre $N_1$ de points fixes de $V$ est divisible par $q$.
\end{surprise}

\subsection{}
Si $r = 2$, la formule des traces est un peu moins effrayante. Drinfeld [Dr] l'a utilis\'ee pour traiter du cas o\`u $S_0$ est vide. Flicker (non publi\'e) a trait\'e du cas o\`u $|S_0| = 1$ et o\`u la monodromie locale en chaque $s \in S$ est unipotente avec un seul bloc de Jordan (c'est \`a dire, puisque $r = 2$, non triviale). Rassemblant les informations ainsi obtenues, on obtient la

\begin{surprise}
Soit $E$ l'ensemble des classes d'isomorphie de $\Ql$-faisceaux lisses irr\'eductibles de rang $2$ sur $X - S$, \`a monodromie locale unipotente (tant non triviale que triviale). Le nombre de points fixes de $V := \Frob^*$ sur $E$ ne d\'epend que de $X_0/\F_q$ et de $|S_0|$.
\end{surprise}

C'est \`a cette surprise que 2.16 fait allusion. Si, comme dans [DF], on exige que la monodromie locale en chaque $s \in S$ soit unipotente avec un seul bloc de Jordan le nombre de points fixes d\'epend de $X_0$, de $S$, et aussi de l'action de Frobenius sur $S$ (vue comme classe de conjugaison dans le groupe sym\'etrique $\mathfrak{S}_{|S_0|}$).

\section{Exemples}

\subsection{}
Prenons $X_0 = \mathbf{P}^1$, $|S_0| = 4$, $r = 2$ et soit $E$ l'ensemble des classes d'isomorphie de $\Ql$-faisceaux lisses irr\'eductibles de rang $2$ sur $X - S$, \`a monodromie locale en chaque $s \in S$ unipotente non triviale. Dans ce cas, le nombre de points fixes des it\'er\'es de la permutation $V = \Frob^*$ de $E$ (notations de 2.1) est donn\'e par
\begin{equation}
N_n = q^n.
\tag{5.1.1}
\end{equation}
Si $|S_0| \geq 2$, (5.1.1) est une application du r\'esultat principal de [DF]. Le curieux r\'esultat suivant ([DF] \S 7) permet d'en d\'eduire le cas o\`u $S_0$ est r\'eduit \`a un point ferm\'e de degr\'e $4$.

\begin{proposition}
Chaque permutation $\sigma$ de type $(2, 2)$ de $S$ se prolonge en une projectivit\'e de $\mathbf{P}^1$. L'action de cette projectivit\'e sur $E$ est triviale.
\end{proposition}

La preuve est par r\'eduction \`a un th\'eor\`eme analogue sur $\C$.

\subsection{}
Il est naturel de compl\'eter l'ensemble $E$ de 5.1 en $\overline{E}$, l'ensemble des classes d'isomorphie \`a semi-simplification pr\`es de $\Ql$-faisceaux lisses de rang $2$ sur $X - S$, \`a monodromie locale unipotente en chaque $s \in S$. Parce que $X_0 = \mathbf{P}^1$ et que $|S_0| = 4$, cette compl\'etion $\overline{E}$ ne diff\`ere de $E$ que par l'adjonction d'un point correspondant au $\Ql$-faisceau constant $\Ql^2$.

On note encore $V$ la permutation de $\overline{E}$ induite par $\Frob^*$. Le nombre $\overline{N}_n$ de points fixes de $V^n$ agissant sur $\overline{E}$ est
\begin{equation}
\overline{N}_n = q^n + 1 \quad \text{(points fixes sur } \overline{E}).
\tag{5.3.1}
\end{equation}

\subsection{}
L'analogue complexe de $\overline{E}$ est l'espace de modules $\overline{M}_\C$ des syst\`emes locaux complexes de rang $2$ sur $\mathbf{P}^1 - S$, pour $|S| = 4$, \`a monodromie locale unipotente. Ces syst\`emes locaux sont pris \`a semi-simplification pr\`es. Si on choisit un point base $o \in \mathbf{P}^1 - S$, le foncteur $F \mapsto F_o$ est une \'equivalence de syst\`emes locaux vers les repr\'esentations de $\pi_1(\mathbf{P}^1 - S, o)$, et on peut prendre comme coordonn\'ees sur $\overline{M}_\C$ les traces des \'el\'ements du $\pi_1$. L'espace $\overline{M}_\C$ est affine, purement de dimension $2$, avec un unique point singulier correspondant au syst\`eme local constant $\C^2$.

La singularit\'e est de type $D_4$. Les nombres de Betti non nuls de $\overline{M}_\C$ sont
\begin{equation}
b_0 = 1, \quad b_2 = 1.
\tag{5.4.1}
\end{equation}

\subsection{}
Je n'ai pas fait une v\'erification compl\`ete, mais la situation semble toute pareille pour $X_0 = \mathbf{P}^1$, $|S_0| = 3$, $r = 3$, monodromie locale unipotente. \`A nouveau, si une des monodromies locales n'est pas \`a un seul bloc de Jordan, le $\Ql$-faisceau est extension it\'er\'ee de syst\`emes locaux constants $\Ql$. L'analogue de 5.3 vaut pour $\sigma$ une permutation cyclique des trois points de $S$. Le nombre de points fixes de $V^n$ est encore $q^n + 1$ (avec ``$1$'' donn\'e par les syst\`emes locaux \`a monodromie locale unipotente r\'eductibles, tous de semi-simplifi\'e un syst\`eme local constant). L'unique point singulier de l'analogue complexe est de type $E_6$.

\subsection{}
Prenons $X_0 = \mathbf{P}^1$, $|S_0| = 4$, $r = 2$ et supposons que comme en 3.4 la monodromie locale impos\'ee soit en chaque $s \in S_0$ donn\'ee par deux caract\`eres distincts $\alpha[s_0'], \alpha''[s_0]$ de $k(s_0)^*$. Supposons v\'erifi\'e (3.2.2) et l'hypoth\`ese de position g\'en\'erale (2.10.3). \`A l'aide de (3.5.1), on obtient

\begin{proposition}
Sous les hypoth\`eses de 5.6, on a
\begin{equation}
N_1 = q + 1 + |S_0(\F_q)|.
\tag{5.7.1}
\end{equation}
\end{proposition}

Arinkin a v\'erifi\'e que pour $X_0 = \mathbf{P}^1$ et $r = 2$, si un fibr\'e $E_0$ de rang $2$ sur $X_0$ avec structure parabolique $\alpha_0$ en $S_0$ est ind\'ecomposable, alors $\End(E_0, \alpha)$ est r\'eduit aux scalaires.

Si on d\'efinit $Z$ comme en 3.6, et qu'on passe au faisceau associ\'e pour la topologie de Zariski, on obtient un foncteur repr\'esentable par un espace alg\'ebrique. Dans le cas $|S_0| = 4$, si on prend le degr\'e $d = 1$, il est repr\'esent\'e par la somme de deux copies de $X_0$, recoll\'ees le long de $X_0 - S_0$. Le nombre de points de ce sch\'ema sur $\F_q$ est donc donn\'e par (5.7.1).

\subsection{}
Passons \`a l'analogue complexe. Sur $\mathbf{P}^1 - S$, avec $|S| = 4$, on consid\`ere l'espace de module des syst\`emes locaux complexes de rang $2$ avec monodromie locale impos\'ee en chaque $s \in S$. En $s \in S$, on demande que la monodromie locale soit conjugu\'ee \`a $A_s = \diag(a_s, b_s)$. On suppose $a_s b_s = 1$ et que si pour chaque $s$, $c_s$ est l'un de $a_s$ ou $b_s$, on n'a jamais $\prod_s c_s = 1$. Posons $A := (A_s)_{s \in S}$ et notons $M(A)$ l'espace de modules, comme en 2.10. Il est lisse purement de dimension $2$. On peut regarder $M(A)$, pour $A$ proche de $1$, comme une d\'eformation de $\overline{M}_\C$ consid\'er\'e en 5.4. Dans cette d\'eformation, l'espace des cycles \'evanescents est de dimension $4$ (puisque la singularit\'e de $\overline{M}_\C$ est de type $D_4$). \`A l'infini, la g\'eom\'etrie ne change pas et on a donc comme nombres de Betti non nuls
\begin{equation}
b_0 = 1, \quad b_2 = 5.
\tag{5.8.1}
\end{equation}

\subsection{}
Prenons $X_0 = \mathbf{P}^1$, $|S_0| = 3$, $r = 3$ et supposons que comme en 3.4 la monodromie locale impos\'ee soit en chaque $s \in S_0$ donn\'ee par trois caract\`eres distincts de $k(s_0)^*$. Supposons v\'erifi\'e (3.2.2) et l'hypoth\`ese de position g\'en\'erale 2.10.3. \`A l'aide de (3.5.1) on obtient ici que
\begin{equation}
N_1 = q + 1 + 2|S_0(\F_q)|.
\tag{5.9.1}
\end{equation}
Le faisceau associ\'e \`a $Z$ comme en 3.6 est en effet, pour $d = 1$, repr\'esent\'e par la somme de trois copies de $X_0$, recoll\'ees le long de $X_0 - S_0$. Si $S_0$ consiste en trois points rationnels sur $\F_q$, (5.9.1) dit que $N_1 = (q+1) + 6$, comme sugg\'er\'e par la singularit\'e $E_6$ de l'analogue complexe, pour une monodromie locale unipotente (cf 5.5).

\section{Rang 1}

\subsection{}
Notons $E(1, \emptyset)$ l'ensemble $E(R)$ de 2.1 pour $r = 1$ et $S_0 = \emptyset$. Le produit tensoriel le munit d'une structure de groupe ab\'elien. La th\'eorie du corps de classe identifie $E(1, \emptyset)^V$ au groupe des caract\`eres \`a valeurs dans $\Ql^*$ du groupe fini $\Pic^0(X_0)(\F_q)$. Cette th\'eorie identifie en effet le groupe fondamental rendu ab\'elien de $X_0$ au compl\'et\'e profini de $\Pic(X_0)$. Le groupe $\Pic(X_0)$ s'envoie sur $\Z$ par l'application degr\'e, avec pour noyau le groupe fini $\Pic^0(X_0)(\F_q)$ des $\F_q$-points de la jacobienne $\Pic^0(X_0)$. Les classes d'isomorphie de $\Ql$-faisceaux $F_0$ lisses de rang un sur $X_0$ s'identifient ainsi aux caract\`eres $\chi$ de $\Pic(X_0)$ \`a valeurs dans $\Zl^*$, la $\F_q$-torsion par $\lambda$ correspond au produit par le caract\`ere $\lambda^n$ de $\Z$, et, attachant \`a $F_0$ la restriction de $\chi$ \`a $\Pic^0(X_0)(\F_q)$, on obtient une bijection entre classes de $\F_q$-torsion de $F_0$ lisses de rang $1$ sur $X_0$ et caract\`eres de $\Pic^0(X_0)(\F_q)$.

\subsection{}
Quels que soient $r \geq 1$ et $R$ comme en 2.1, le produit tensoriel induit une action du groupe ab\'elien $E(1, \emptyset)$ sur $E(R)$, et de $E(1, \emptyset)^V$ sur $E(R)^V$. Si $r = 1$, et que $R$ v\'erifie la compatibilit\'e (2.10.1), le corps de classe montre que $E(R)^V$ est un espace principal homog\`ene sous $E(1, \emptyset)^V$. On a donc
\begin{equation}
|E(R)^V| = |E(1, \emptyset)^V| = |\Pic^0(X_0)(\F_q)| \quad \text{(si } r = 1 \text{ et (2.10.1))}.
\tag{6.2.1}
\end{equation}

Si $r > 1$, prendre la puissance ext\'erieure $\bigwedge^r$ fournit une application
\[
\det: E(R) \longrightarrow E(\textstyle\bigwedge^r R)
\]
telle que, pour l'action, not\'ee $\otimes$, de $E(1, \emptyset)$, on ait
\[
\det(\mathcal{L} \otimes F) = \mathcal{L}^r \otimes \det(F).
\]
De m\^eme, apr\`es passage aux points fixes de $V = \Frob^*$, pour $\det: E(R)^V \to E(\bigwedge^r R)^V$ et l'action de $E(1, \emptyset)^V$.

L'action de $E(1, \emptyset)^V$ sur $E(R)^V$ n'est pas n\'ecessairement libre, et les fibres de
\[
\det: E(R)^V \longrightarrow E(\textstyle\bigwedge^r R)^V
\]
n'ont pas n\'ecessairement toutes le m\^eme nombre d'\'el\'ements. N\'eanmoins, dans tous les cas qui ont pu \^etre calcul\'es $|E(1, \emptyset)^V|$ divise $|E(R)^V|$. Comme fonction de $n$, $|\Pic^0(X_0)(\F_{q^n})|$ est de la forme (2.15.1):
\[
|\Pic^0(X_0)(\F_{q^n})| = \sum_i (-1)^i \Tr(\Frob^*, H^i(\Pic^0(X), \Q_{\ell'}))
\]
\[
\phantom{|\Pic^0(X_0)(\F_{q^n})|} = \sum_i (-1)^i \Tr(\textstyle\bigwedge^i \Frob^*, H^1(X, \Q_{\ell'})).
\]

Je conjecture la forme pr\'ecis\'ee suivante de la conjecture 2.15 (i)

\begin{conjecture}
Comme fonction de $n$, le nombre de points fixes $N_n(R)$ de $V^n$ est de la forme
\begin{equation}
N_n(R) = |\Pic^0(X_0)(\F_{q^n})| \cdot \sum_k c_k \gamma_k^n
\tag{6.3.1}
\end{equation}
pour des entiers $c_k$ et des nombres $\gamma_k$ convenables.
\end{conjecture}

\subsection{}
Voici un analogue de 6.3 en g\'eom\'etrie alg\'ebrique. Soient $A_0$ une vari\'et\'e ab\'elienne sur $\F_q$ et $f_0: X_0 \to A_0$ un morphisme propre et lisse. Bien que les fibres de $f: X_0(\F_{q^n}) \to A_0(\F_{q^n})$ n'aient pas n\'ecessairement le m\^eme nombre d'\'el\'ements (exemple: $X_0 = A_0$ et $f = \text{multiplication par un entier premier \`a } p$), $|A_0(\F_{q^n})|$ divise $|X_0(\F_{q^n})|$. En effet, les $R^i f_! \Q_\ell$ sont des $\Q_\ell$-faisceaux lisses sur $A$, et le groupe fondamental (ab\'elien) de $A$ agit sur ces faisceaux. On sait que cette action est semi-simple. Soit $(R^i f_! \Q_\ell)^0$ le sous-faisceau des invariants. Il provient d'un sous-faisceau $(R^i f_{0!} \Q_\ell)^0$ de $R^i f_{0!} \Q_\ell$ sur $A_0$, et ce dernier est l'image inverse sur $A_0$ d'un faisceau sur $\Spec(\F_q)$, que nous noterons $H^i(X_0/A_0)$. On sait que
\begin{equation}
H^i(A, (R^j f_! \Q_\ell)^0) \xrightarrow{\sim} H^i(A, R^j f_! \Q_\ell).
\tag{6.4.1}
\end{equation}
La formule des traces de Grothendieck pour le nombre de $\F_q$-points de $X_0$ peut donc se r\'e\'ecrire
\begin{equation}
|X_0(\F_{q^n})| = \Tr\Bigl(\Frob^*, \sum_i (-1)^i H^i(A, \Q_{\ell'}) \otimes \sum_j (-1)^j H^j(X_0/A_0)\Bigr).
\tag{6.4.2}
\end{equation}
Si les $\gamma_k$ sont les valeurs propres de Frobenius sur les $H^j(X_0/A_0)$, et les $c_k$ la somme altern\'ee de leurs multiplicit\'es, on d\'eduit de (6.4.2) que
\begin{equation}
|X_0(\F_{q^n})| = |A_0(\F_{q^n})| \cdot \sum_k c_k \gamma_k^n.
\tag{6.4.3}
\end{equation}
La divisibilit\'e promise de $|X_0(\F_{q^n})|$ par $|A_0(\F_{q^n})|$ r\'esulte d'une formule analogue \`a (6.3.1). Les valeurs propres de Frobenius sur $H^j(X_0/A_0)$, contenu dans la cohomologie d'une fibre de $X_0 \to A_0$, sont des entiers alg\'ebriques.

\subsection{}
Dans 6.4, le point crucial n'est pas que $f_0$ soit propre et lisse, mais que les $R^i f_{0!} \Q_\ell$ soient lisses. On peut \'eviter d'avoir \`a supposer que les $R^i f_{0!} \Q_\ell$ soient semi-simples en consid\'erant, plut\^ot que le sous-faisceau des invariants pour l'action de $\pi_1(A)$, le plus grand sous-faisceau sur lequel l'action est unipotente.

Supposons qu'il existe une action de $A_0$ sur $X_0$, et un entier $r$, tels que $f_0$ soit \'equivariant pour l'action de $A_0$ sur $A_0$ par $x: y \mapsto rx + y$. Sur cette hypoth\`ese, les images inverses des $R^i f_{0!} \Q_\ell$ sur $A$ par la multiplication par $r: A \to A$, sont des faisceaux constants. Ceci implique et leur lissit\'e, et leur semi-simplicit\'e.

\subsection{}
Si $g \geq 1$, la conjecture 6.3 implique que dans l'expression conjecturale (2.15.1) pour $N_n(R)$, on a $\sum_i a_i = 0$.

Passons aux analogues complexes comme en 2.10. Soit $M(1, \emptyset)$ le groupe des classes d'isomorphie de syst\`eme locaux complexes de rang $1$ sur $\mathcal{P}$. Il est isomorphe \`a $\C^*{}^{2g}$ et, si la compatibilit\'e $\det R^*(s) = 1$ est v\'erif\'ee, $M(\det R^*)$ est un espace principal homog\`ene sous $M(1, \emptyset)$. Le groupe $M(1, \emptyset)$ agit sur $M(R^*)$ et sur $M(\det R^*)$. Notons $\otimes$ l'action. L'application induite par $\bigwedge^r$:
\[
\det: M(R^*) \longrightarrow M(\det R^*)
\]
v\'erifie $\det(\mathcal{L} \otimes F) = \mathcal{L}^r \otimes \det(F)$. Ceci implique que $\det$ est une fibration, que les $R^i \det_! \Z$ sont localement constants, et que l'action de $\pi_1(M(\det R^*))$ sur les $R^i \det_! \Z$ se factorise par son quotient $\pi_1/r\pi_1 \cong (\Z/r\Z)^{2g}$.

Si $g \geq 1$, on en d\'eduit, par des arguments analogues \`a ceux de 6.4, que $\chi(M(R^*)) = 0$, en accord avec 2.15 (ii). Soit $(R^i \det_! \Q)^0$ le sous-syst\`eme local de $R^i \det_! \Q$ des invariants sous l'action de $\pi_1$, et $\chi(M(R^*)/M(\det R^*))$ la somme altern\'ee des rang des $(R^i \det_! \Q)^0$.

\begin{conjecture}
Avec les notations de 6.3 et de 6.6, on a
\[
c_k = \chi(M(R^*)/M(\det R^*)).
\]
\end{conjecture}

\section*{Bibliographie}

\begin{enumerate}
\item[{[DF]}] P. Deligne and Y. Z. Flicker. Counting local systems with principal unipotent local monodromy, \textit{Annals of Math.} (to appear).
\item[{[Dr]}] V. Drinfeld. The number of two dimensional irreducible representations of the fundamental group of a curve over a finite field, \textit{Funk. anal. i Prilozhen.} 15 4 (1981), 75--76.
\item[{[L]}] L. Lafforgue. Chtoucas de Drinfeld et correspondance de Langlands, \textit{Inv. math.} 147 1 (2002), 1--241.
\end{enumerate}

\end{document}
